% A PARAMETER IS WRITTEN OUT WITH THE CHARACTER ITS PARAMETER TEXT USED.
%
% More than one character can have category six at once, and a macro
% defined with the second one shows its parameters with THAT character:
%
%   \catcode`\!=6  \def\e!1\e{!1!1}      \show\e  gives  !1\e ->!1!1
%
% The rule is one character kept while a token list is written out. It
% starts as `#' for each list and is set by every parameter of the
% MACRO'S PARAMETER TEXT as that goes past; the parameters of the body
% are written with whatever it was left at. So a macro whose parameter
% text uses both characters shows its whole body with the LAST of them:
%
%   \def\f#1!2{[#1][!2]}   \show\f  gives  #1!2->[!1][!2]
%
% and the `#1<-' that \tracingmacros puts in front of an argument uses
% the character the macro's own parameter text left behind, because the
% two are written one after the other.
%
% HSTeX stored only the parameter's NUMBER, so every parameter came out
% as `#' whatever it was written with.
\catcode`\{=1 \catcode`\}=2 \catcode`\#=6
\tracingonline=1
\catcode`\!=6
\def\d#1\d{#1#1}
\def\e!1\e{!1!1}
\def\f#1!2{[#1][!2]}
\message{[A a macro written with hash]}
\show\d
\message{[B a macro written with the other parameter character]}
\show\e
\message{[C both in one]}
\show\f
\message{[D and how they trace]}
\tracingmacros=2
\message{\d X\d}
\message{\e Y\e}
\tracingmacros=0
\message{[done]}
\end
